\scanpage{7}
\margin{à vérifier}
Il revient au $=$ de se donner une application
\[
\unsure{\varphi|\phi}_R:\underline{R}\longrightarrow\deleted{\unsure{\operatorname{Repaff}}(E)}\operatorname{Drap.aff}_n(E)
\]
\margin{drapeaux maximaux de $E$}
et \deleted{un homo} et un \deleted{tel qu'il existe un unique} homo
\diagram{\begin{tikzcd}[ampersand replacement=\&]
G \arrow[r,"\varphi_G"] \& \operatorname{Aut}\operatorname{aff}E
\end{tikzcd}}
compatible avec $\varphi_R$, tel que les sommets des \deleted{repères} drapeaux de l'image de $\underline{R}$ engendrent affinement $E$. (ce qui implique l'unicité de $\varphi_G$ \unsure{élément} par \unsure{élément} \dots), et la condition suivante :

\inserted{le sous-groupe $G_{f_i}$ de $G$ stabilisateur de la composante $f_i$ de dim $i$ d'un ``repère'' $\underline{r}\in\underline{R}$, opérant sur $E$, \unsure{invarie} \deleted{la composante} $F_i$ de l'image $r$ de $\underline{r}$ par $\varphi$}
% insertion continues into the left margin, slanted

Comme $\underline{R}$ est un \unsure{torseur|espace} sous $G$, il suffit en fait, étant choisi une origine $\underline{r}\in\underline{R}$, \inserted{permettant d'identifier $\underline{R}\unsure{:|\simeq}R_{p_*}$}, de se donner $\varphi_{G\inserted{=\Gamma}}$ et $\varphi_R(\underline{r})=r\in\operatorname{Drap.aff}_n(E)$, avec comme seules conditions que

a) \deleted{$\unsure{\{}$} Si $j\neq i$, $\sigma_j$ \unsure{invarie} $f_j$

b) les $g.f_0$ ($g\in G=\Gamma$) engendrent \inserted{affinement $E$ fibre par fibre \dots}
% archive sheet number: 6
